\section{Finite-range integration and local partition algebra}\label{rg:section}

The common task is to turn a local activity expansion into an exact
expansion on coarser blocks, with bounds uniform in volume.  The later
spin and height sections supply their own initial activities to these
operations.

This section fixes the analytic conventions used by the spin and height constructions.  A field is real before analytic continuation.  Counting measure is used on vertices and on unoriented edges.  On the nearest-neighbor lattice let $\partial$ be the edge gradient and $\mathcal L=\partial^*\partial$.  Constants are averaged uniformly modulo a specified period $\omega$, separately from the centered Gaussian integral.  Consequently an independent Gaussian constant can be added at any integration step without changing a periodic partition integral.  Bounds below are uniform in the volume.  Constants can depend on fixed geometric parameters, but a constant asserted independent of $L$ is chosen before the block ratio $L$.

Finite-range Gaussian decompositions were developed by
Brydges, Guadagni, and Mitter \citep{BrydgesGuadagniMitter2004}
and, for Green forms under a finite-propagation hypothesis, by Bauerschmidt
\citep{Bauerschmidt2013}.  The analytic polymer and extraction
framework follows the line of Brydges--Yau and subsequent
renormalisation constructions \citep{BrydgesYau1990,Brydges2009,BPR2024I}.
We give the polynomial kernels and every estimate needed here,
including the boundary versions, analytic norms, and regulator reserve.

\subsection{Polynomial covariances and their boundary versions}

\begin{lemma}[Finite-range decomposition]\label{rg:kernels}
Set $p=32$, $\gamma=1/16$, and assume $|s|\le1/4$.  For integral $t\ge1$ define the polynomial
\[
 Q_t(\cos z)=\left(\frac{\sin(tz/2)}{t\sin(z/2)}\right)^{2p}.
\]
Let $L$ be a power of two, $S=Lm$, and $m$ a power of $L$.  The operator with multiplier
\begin{equation}\label{rg:shell}
 \Gamma_m(\lambda)=\frac{1-\gamma\lambda}{(1+s)\lambda}
 \{Q_m(1-\lambda/16)-Q_S(1-\lambda/16)\}
\end{equation}
is positive semidefinite on the plane, the square torus, and the free square of cells.  It has range at most $2pS$.  The quotient at zero is its polynomial continuation.  Together with independent variance $\gamma/(1+s)$ and the terminal covariance
\begin{equation}\label{rg:terminal-kernel}
 \Gamma_*(\lambda)=\frac{(1-\gamma\lambda)Q_{m_*}(1-\lambda/16)}{(1+s)\lambda},
 \qquad \lambda>0,
\end{equation}
the shells with upper scale at most $m_*$ sum to $(1+s)^{-1}\mathcal L^{-1}$ on nonconstant modes.  Here $D_0\ge2$ is fixed sufficiently large.  In a square of side $M>D_0$, let $m_*$ be the largest nonnegative power of $L$ strictly below $M/D_0$.  For $M\le D_0$, set $m_*=1$ and use no shells.  Since $Q_1=1$, the latter convention gives the same decomposition; for $M=1$ there are no nonconstant modes.  The constant mode of the terminal covariance can be set to zero.

For the plane shells and the finite-volume shells included above,
\begin{align}
 |\Gamma_m(x,y)|&\le C\left(1+\log^+\frac{S}{m+\dist(x,y)}\right),
 &\Gamma_m(x,x)&\ge c\log L,\label{rg:kernel-pointwise}\\
 \|\Gamma_m\|_{2\to2}&\le CS^2,
 &\|\nabla^r\Gamma_m^{1/2}\|_{2\to2}^2&\le C_r m^{-2(r-1)},\quad1\le r\le6.\label{rg:kernel-operator}
\end{align}
For mixed differences in the two kernel arguments of total order
$r=|a|+|c|$, with the reflected difference conventions at a free side,
\[
 |\nabla_x^a\nabla_y^c\Gamma_m(x,y)|
 \le C_r\bigl(m+\dist(x,y)\bigr)^{-r},
 \qquad 1\le r\le12.
\]
In particular $\E|\nabla^r\zeta(x)|^2\le C_r m^{-2r}$ for $1\le r\le6$ when $\zeta$ has covariance $\Gamma_m$.  These positive-order diagonal and operator estimates also hold for $\Gamma_*$ at scale $m_*$, with constants independent of $D_0,L$ when $M\ge2m_*$, as the stopping convention guarantees for $M\ge2$.  For $M=1$ all positive differences vanish.
\end{lemma}

\begin{proof}
The square of the Dirichlet sum is a polynomial in $\cos z$ of degree $t-1$, so $Q_t$ has degree $p(t-1)$.  For $t$ integral,
\[
 Q_{2t}(\cos z)=Q_t(\cos z)\cos^{2p}(tz/2)\le Q_t(\cos z).
\]
Iteration proves positivity of every difference in \eqref{rg:shell}.  Since $Q_t(1)=1$, division by $\lambda$ is removable, and multiplication by $1-\gamma\lambda$ leaves degree at most $p(S-1)$.  The spectrum of $\mathcal L$ lies in $[0,8]$, where $1-\gamma\lambda\ge1/2$.  A polynomial of degree $d$ in a nearest-neighbor operator has range at most $d$.  This proves positivity and the stated, slightly enlarged, range.  Telescoping proves the decomposition, including the empty shell sum when $m_*=1$ because $Q_1=1$.  Its continuation at zero is positive since
\[
 Q_t(1-\lambda/16)=1-\frac{t^2-1}{3}\lambda+O(\lambda^2).
\]

Write a shell as the sum of its binary pieces $H_t$, $t=m,2m,\ldots,S/2$.  On the Fourier square the nearest-neighbor symbol is
$\lambda(\xi)=4\sum_{i=1}^2\sin^2(\xi_i/2)\asymp |\xi|^2$.
For $0\le\lambda\le8$, the identity just used and $|\sin u|\le\min(1,|u|)$ give
\begin{equation}\label{rg:binary-symbol}
 0\le H_t(\lambda)\le C\min\{t^2,\lambda^{-1}\min(1,(t\sqrt\lambda)^{-2p})\}.
\end{equation}
Indeed the difference factor $1-\cos^{2p}(tz/2)$ is at most $Ct^2z^2$, whereas $Q_t(\cos z)\le C\min(1,(tz)^{-2p})$; here $z\asymp\sqrt\lambda$.  A difference of order $r$ has Fourier multiplier at most $C_r|\xi|^r$.  Integrating \eqref{rg:binary-symbol} in polar coordinates, split at $|\xi|=t^{-1}$, gives
\[
 \int |\xi|^r H_t(\lambda(\xi))\,\frac{\dd\xi}{(2\pi)^2}\le C_rt^{-r},\qquad0\le r\le12.
\]
The corresponding normalized Fourier sums obey the same inequality whenever $t\le M/D_0$: divide the frequency square into annuli of width $M^{-1}$ and use that their number of points is bounded by a constant times their area times $M^2$, with the zero-frequency contribution $O(t^2/M^2)$.  Taking suprema instead of integrals gives
$\sup_\xi |\xi|^{2r}H_t\le C_rt^{2-2r}$.
For the full shell and $r=1$, use directly $\lambda\Gamma_m\le C$; for $r>1$ the binary bounds form a geometric series.  For $r=0$ their operator bounds sum to $CS^2$.  Their kernels vanish beyond distance $2pt+O(r)$, so only $t\ge c\dist(x,y)$ can contribute.  Counting these dyadic scales proves the logarithmic bound and summing $t^{-r}$ proves the positive-order bounds.

Even reflection of the cell coordinates across $-1/2$ and $M-1/2$ identifies the free operator with the even subspace of the $2M$ torus operator.  Its kernel is the sum of the four reflected images in two dimensions.  Since the range is less than a fixed small fraction of $M$, at most a fixed number of periodic images enter the preceding estimates.  The reflected difference conventions are the ones in those estimates.  Orthogonal restriction, with the reflection multiplicity in the counting norm, also proves the operator estimates.

For completeness, the lower bound is not inferred from positivity of image terms.  Choose a small fixed $a>0$.  Whenever both coordinates of the Fourier frequency are between $a/t$ and $4a/t$, the binary multiplier is at least $c_at^2$.  This follows by the expansion at zero of $Q_t-Q_{2t}$, uniformly on the indicated compact annulus after decreasing $a$.  In the plane the annulus has area at least $c_at^{-2}$.  On the torus it has at least $c_a(M/t)^2$ points for $D_0$ sufficiently large.  In the free cosine basis the factors at cell coordinate $i$ are $\cos(\pi u(i+1/2)/M)$.  Pair each index $u$ in a subinterval $[aM/t,2aM/t]$ with $2u$.  The elementary inequality
\[
 \cos^2\theta+\cos^2(2\theta)\ge7/16
\]
ensures a uniform lower bound on each pair, and each frequency has bounded pairing multiplicity.  Taking the product of the two one-dimensional sums proves $H_t(x,x)\ge c$ also at a corner.  There are $\log_2L$ binary pieces, proving the diagonal assertion.

For the terminal remainder use $\lambda^{-1}Q_{m_*}\le C\lambda^{-1}\min\{1,(m_*\sqrt\lambda)^{-2p}\}$.  A positive-order diagonal integral is finite at zero and bounded by $C_rm_*^{-2r}$; the discrete sum has the same bound uniformly for $M\ge2m_*$: there are at most $C(M/m_*)^2$ modes in the disk $|\xi|\le m_*^{-1}$, and the remaining annuli form a convergent power sum.  The operator bounds follow from the same multiplier with $\lambda\ge cM^{-2}$.  Reflection gives the free version.  For $M=1$ there are no nonconstant modes.  No estimate on the value of the terminal constant mode is needed.
\end{proof}

\begin{lemma}[Killed Green function]\label{rg:green}
For the nearest-neighbor killed Laplacian of the spin square, let $G_n$ be its inverse and put $M=2n$.  In lattice units, for interior vertices and differences interpreted by odd reflection across the zero sides,
\begin{align}
 0\le G_n(x,y)&\le C\left(1+\log^+\frac{M}{1+|x-y|}\right),\label{rg:green-value}\\
 |\nabla_x^rG_n(x,y)|&\le C_r(1+|x-y|)^{-r},\qquad1\le r\le6.\label{rg:green-difference}
\end{align}
After identifying vertices with their positions in $(-1,1)^2$, $G_n(x,y)$ converges locally uniformly off the diagonal to the continuum Dirichlet Green kernel.  Positive-order differences, observed in lattice coordinates a bounded distance from a source whose distance from the walls tends to infinity, converge to their infinite-lattice values.  These conclusions apply to rotated gradients and to their local stream lifts, with even reflection of the lift at a free side.
\end{lemma}

\begin{proof}
Use the decomposition of Lemma~\ref{rg:kernels} with the factors $(1-\gamma\lambda)/(1+s)$ omitted.  Odd reflection gives its Dirichlet version.  Sum binary pieces up to a scale comparable to $M$ and keep the smoothed positive-spectrum remainder.  A piece of scale $t$ has differences bounded by $C_rt^{-r}$ and vanishes when $t<c|x-y|$, apart from reflected images whose distances are no smaller than the direct distance.  Summing proves \eqref{rg:green-value} and \eqref{rg:green-difference}; nonnegativity follows independently from the maximum principle.  The terminal remainder has differences $O(M^{-r})$ and value $O(1)$, since the smallest Dirichlet eigenvalue is at least $cM^{-2}$.

For the local lattice limit, all fixed-scale pieces eventually coincide with the plane pieces by finite range.  The tail of a derivative of positive order is bounded by $C\sum_{t\ge T}t^{-r}\le C_rT^{-r}$.  Taking the volume limit and then $T\to\infty$ proves convergence.  For the continuum limit, the discrete sine modes converge to the continuum sine modes: on every fixed mode $n^2\lambda_n$ converges to its continuum eigenvalue.  If $f,g$ are smooth and compactly supported in the interior, repeated summation by parts against those sine modes bounds their discrete coefficients by $C_q(1+|k|)^{-q}$, uniformly in $n$, for every fixed $q$.  The eigenvalue inequality $n^2\lambda_n(k)\ge c|k|^2$ then permits dominated summation, identifying the distributional limit of the kernels with $(-\Delta_D)^{-1}$.  Equation~\eqref{rg:green-difference}, multiplied by $n$ for first differences on compact subsets away from the diagonal, gives equicontinuity; \eqref{rg:green-value} gives boundedness there.  Distributional identification therefore upgrades to local uniform convergence.  Rotating an odd-reflected potential gradient makes its normal component odd and its tangential component even.  Its curl vanishes away from source plaquettes; path integration there supplies an even-reflected local lift with the asserted difference estimates.
\end{proof}

An activity records a local correction to a Gaussian reference on a
connected collection of blocks.
The analytic norms control field derivatives, while the Gaussian
regulator permits controlled growth at large fields.
The localization bounds treat neutral, charged, and restricted terms
separately, and Proposition~\ref{rg:partition-map} combines the
resulting local operations into an exact partition map.

\subsection{Blocks, analytic norms, and the Gaussian regulator}

Fix $r_0=4p+20$ and $R=r_0+8$, increasing both when a model has a larger finite interaction range.  A polymer is a finite set of blocks connected by steps of sup block-index length at most $R$.  Its size $|X|$ is its number of blocks.  At scale $m$ the plane blocks have side $m$.  In finite volume the grid levels are the retained activity levels of the construction in use, each with $m=1$ or $m<M/D_0$ for its fixed sufficiently large clearance $D_0$.  A coarse grid at $S=Lm$ is present only when $S$ is retained.  On a torus of side $M$, require $m\mid M$ at every retained level and use full blocks of side $m$ with one common origin for the nested grids.  This divisibility is a condition for the block calculus; the spectral decomposition of Lemma~\ref{rg:kernels} does not require it.  For free cells $[0,M)^2$, choose one fixed integral translate of all the nested grids.  On each axis retain the cuts in that translate of $m\Z$ whose distance from both endpoints is at least $m$.  At $m=1$ this gives unit intervals.  At every other retained level, the clearance ensures that end intervals have lengths in $[m,2m)$ and other intervals have length $m$.  Coarser cuts are fine cuts, so this defines nested partitions.  Distances on a torus are periodic.  For a retained pair $m,S$, the closure $\bar X$ is the union of the coarse blocks meeting $X$.

\begin{lemma}[Geometry and counting]\label{rg:geometry}
There are constants $C_0,n_0$ and a lower bound for $L$, depending only on the fixed connection radii, such that
\begin{equation}\label{rg:closure}
 |\bar X|\le C_0(1+R|X|/L),\qquad |\bar X|\le |X|,
\end{equation}
$\bar X$ is $r_0$-connected, and $|X|>n_0$ implies $|\bar X|\le|X|/2$.  If $|X|\le n_0$, the index diameter of $\bar X$ is at most one.  The number of size-$n$ fine polymers containing a fixed fine block is at most $C_R^n$.  For a fixed coarse root, configurations of $j$ mutually compatible fine polymers, with total fine size $n$ and $r_0$-connected union of closures containing that root, number at most
\begin{equation}\label{rg:family-count}
 C^n(CL^2)^{2j}.
\end{equation}
These statements hold for the irregular end blocks with the same constants.
\end{lemma}

\begin{proof}
Choose a spanning tree of the fine block graph of $X$ and traverse it in both directions.  Its traversal has fewer than $2|X|$ steps, each of spatial length at most $C Rm$.  Break it after each accumulated length $S/4$.  A resulting piece has diameter at most $S/2$ and meets a bounded number of coarse rectangles because every coarse side is at least $S$.  This proves the first inequality in \eqref{rg:closure}; the second follows from nesting.  Successive fine blocks in the traversal meet coarse blocks at index distance at most $1+CR/L$, which is less than $r_0$ for large $L$.  Take $n_0>4C_0$ and then $L>4C_0R$.  The half-size assertion follows.  Increase $L$ further so $CRn_0m<S/2$.  A coordinate interval of this length cannot meet three consecutive coarse intervals, proving the diameter assertion.

A canonical depth-first traversal records a rooted animal by a word of length at most $2n$ in a fixed alphabet of allowed fine steps, giving the first count.  To encode a family, join its $j$ fine components by a tree using pairs of blocks whose containing coarse blocks have distance at most $r_0$.  Each joining displacement has length at most $CL$ in fine block indices and hence has at most $CL^2$ choices.  Traverse fine component trees and the joining tree together.  A word specifying step types, component boundaries, and traversal returns has at most $C^n$ choices; the long steps are traversed at most $2(j-1)$ times.  The first fine root has at most $CL^2$ positions in the specified coarse block.  This gives \eqref{rg:family-count}, with one unused factor.  All estimates use only the lower and upper bounds $m,2m$ on side lengths and therefore include end blocks.  In the finite-volume decomposition above, whenever a shell is used its upper scale is at most $m_*<M/D_0$ with $D_0$ larger than the fixed neighborhood radii, so every small neighborhood used here unwraps on a torus.
\end{proof}

Let $D_m(X)$ consist of the sites in $X$ and all blocks at index distance at most two.  Activities have field support in $X$ with one lattice step of collar; a fixed larger microscopic collar is permissible if it is at most a fixed small fraction of $m$.  Write $D=D_m(X)$ and define
\begin{align}
 g_D(\phi)&=\frac12\sum_{v\in D}\sum_{e=\pm e_i}|\nabla_e\phi(v)|^2,\label{rg:reg-def}\\
 P_m(D,\phi)&=m\sum_{\{v,w\}\in\partial D}|\phi(v)-\phi(w)|^2,\nonumber\\
 T_m(D,\phi)&=\sum_{B\subset D}m^4\max_{v\in B,\,|a|=2}|\nabla^a\phi(v)|^2,\nonumber\\
 U_m(X,\phi)&=g_D(\phi)+\delta P_m(D,\phi)+T_m(D,\phi),
 \qquad W_m^\kappa(X,\phi)=e^{\kappa U_m(X,\phi)}.\nonumber
\end{align}
Here $\partial D$ consists of edges with exactly one endpoint in $D$.  Difference stencils based in a block may leave it; they remain inside its fixed stencil neighborhood.  At a free wall the field is extended evenly, so the crossing first difference is zero.

For a direction $f$ put
\[
 \|f\|_{m,D}=\max_{v\in D}\{|f(v)|,m|\nabla f(v)|,m^2|\nabla^2f(v)|\}.
\]
With $h_0$ fixed, define the analytic point norm and activity norm by
\begin{equation}\label{rg:analytic-norm}
 \begin{aligned}
 |F|_{m,X,\phi}&=\sum_{a\ge0}\frac{h_0^a}{a!}
 \sup_{\|f_i\|_{m,D_m(X)}\le1}|D^aF(\phi)[f_1,\ldots,f_a]|,\\
 \|F\|_{m,A,\kappa}&=\sup_{X,\phi}
 \frac{A^{|X|}|F(X)|_{m,X,\phi}}{W_m^\kappa(X,\phi)}.
 \end{aligned}
\end{equation}
Derivatives are real directional derivatives with locally convergent analytic extensions.  The derivative sum bounds translation by a complex direction of norm less than $h_0$.  Leibniz's formula gives the point-norm product inequality, since the binomial coefficient cancels the factorials.  Restriction of a coarse test direction to an included fine collar has norm at most one.

\begin{lemma}[Regulator stability with an unchanged bulk coefficient]\label{rg:regulator}
There are choices in the order
\begin{equation}\label{rg:choice-order}
 \delta>0\ \hbox{small},\qquad \kappa>0\ \hbox{small depending on }\delta,
 \qquad L\ \hbox{large},
\end{equation}
independent of the volume, for which
\begin{equation}\label{rg:reg-stability}
 \E W_m^\kappa(X,\phi+\zeta)\le2^{|X|}W_S^\kappa(Y,\phi),\qquad \bar X\subset Y,
\end{equation}
when $\zeta$ has covariance $\Gamma_m$.  Connectivity is unnecessary.  The estimate holds also with $2\kappa$ in place of $\kappa$.  For a connected $X$ with $|X|\le n_0$, the right side may instead be $2^{|X|}W_S^{\kappa/2}(\bar X,\phi)$.  At a fixed scale $X\subset X'$ implies $U_m(X,\phi)\le U_m(X',\phi)$.  Products on families whose collars are separated beyond the fluctuation range satisfy the corresponding estimate for the union.  At the terminal scale,
\begin{equation}\label{rg:terminal-reg}
 \E W_{m_*}^\kappa(X,\zeta_*)\le C^{|X|}
\end{equation}
for sufficiently small $\kappa$, with fixed-volume-independent $C$.
\end{lemma}

\begin{proof}
We first record the elementary estimates underlying the claim.  If $B$ is any rectangle of side lengths in $[m,2m]$ and $a$ is a first difference based on its stencil neighborhood, averaging along coordinate paths gives
\begin{align}
 m\sum_{v\in F}|a(v)|^2&\le C\sum_{v\in B}|a(v)|^2+Cm^4\max_B|\nabla a|^2,\label{rg:trace}\\
 m^2\max_B|a|^2&\le C\sum_{v\in B}|a(v)|^2+Cm^4\max_B|\nabla a|^2,\label{rg:gradient-sup}
\end{align}
where $F$ is one face, including an adjacent stencil face.  To verify these inequalities, compare $a(v)$ with $a(w)$ along a path of at most $4m+4$ unit steps; its difference is at most $Cm\max|\nabla a|$.  Square the comparison and average over all $w\in B$.  For \eqref{rg:trace}, sum over its at most $2m$ face sites and multiply by $m$.  All reflected versions follow by applying the same comparison to the even extension.

Put $D'=D_S(Y)$ and $E=D'\setminus D$.  The collars and nesting imply $D\subset D'$.  Counting at most $CL^2$ fine blocks in each coarse block gives
\begin{align}
 T_m(D)&\le CL^{-2}T_S(D'),\label{rg:det-comparison}\\
 P_m(D)&\le L^{-1}P_S(D')+Cg_E+CL^{-2}T_S(D').\nonumber
\end{align}
For the second inequality an old boundary face that survives on $\partial D'$ keeps the same edge increments with coefficient $m/S$.  Every other old face abuts a fine block in $E$.  Apply \eqref{rg:trace} to that block, basing a crossing increment at its outside endpoint.  Each block is charged on at most four faces.  The first term is bounded by $Cg_E$, and the second sums as in the first line.  The same argument for two sets at the same scale charges disappearing faces to the added blocks; their $g$ and $T$ increments absorb $\delta$ times this cost when $\delta C\le1$.  This proves monotonicity, including a nonconnected added region.

The pure-fluctuation estimate is uniform in $L$.  The one-dimensional inequality
\[
 |u(i)|\le\frac1{q}\sum_{j=1}^{q}|u(j)|+\sum_{j=1}^{q-1}|u(j+1)-u(j)|
\]
on an interval of length $q\asymp m$, applied in each coordinate and followed by Cauchy--Schwarz, yields the two-dimensional discrete Sobolev bound
\[
 \max_B |u|^2\le C\sum_{a\in\{0,1\}^2}m^{2|a|-2}\sum_{v\in B^+}|\nabla^a u(v)|^2.
\]
Here $B^+$ adds only the fixed difference stencils and has bounded overlap as $B$ varies.  Apply it to $u=m\nabla\zeta$ and $u=m^2\nabla^2\zeta$.  It bounds $U_m(X,\zeta)$ by a fixed multiple of the squared Euclidean norm of the stacked Gaussian vector
\[
 Z=(m^{j-1}\nabla^j\zeta(v):1\le j\le4,\ v\in D^+).
\]
The operator norm of its covariance is at most $C$ by \eqref{rg:kernel-operator}; its trace is at most $C|X|$ because each diagonal entry is $O(m^{-2})$ and $|D^+|\le C m^2|X|$.  If its eigenvalues are $\lambda_i$, then
\[
 \E e^{a|Z|^2}=\prod_i(1-2a\lambda_i)^{-1/2}
 \le\exp\{2a\textstyle\sum_i\lambda_i\},\qquad 2a\max_i\lambda_i\le1/2.
\]
Thus $\E e^{aU_m(X,\zeta)}\le e^{Ca|X|}$ for all sufficiently small $a$.  Its exponential cost can be made less than $2^{c|X|}$ with any fixed $c>0$.  The terminal positive-difference bounds prove \eqref{rg:terminal-reg} by this same calculation.

It remains essential to treat the bulk cross term without a squared triangle inequality.  For an edge $e=\{v,w\}$ set $a_H(e)=(\ind_H(v)+\ind_H(w))/2$, and put
\[
 B_H(\phi,\zeta)=\sum_ea_H(e)\partial\phi(e)\partial\zeta(e),\qquad
 g_H(\phi+\zeta)=g_H(\phi)+2B_H(\phi,\zeta)+g_H(\zeta).
\]
The gradient operator bound gives $\Var B_E\le Cg_E(\phi)$.  Summation by parts on $D'$, orienting crossing edges from $v\in D'$ to $w\notin D'$ and setting $h_e=\phi(v)-\phi(w)$, gives the exact identity
\begin{equation}\label{rg:ibp}
 \partial^*a_{D'}\partial\phi
 =\ind_{D'}\mathcal L\phi-\frac12\sum_{e\in\partial D'}h_e(\delta_v+\delta_w).
\end{equation}
For the first term, $\|\Gamma_m\|\le CS^2$ and the pointwise second-difference bound show
\[
 \Var\ip{\ind_{D'}\mathcal L\phi}{\zeta}
 \le CS^2\sum_{D'}|\mathcal L\phi|^2\le CT_S(D',\phi).
\]
For the boundary term, the absolute covariance row sum over all crossing endpoints is at most $CS$.  Indeed finite range leaves only a bounded number of nearby coarse faces, each of length $O(S)$, and on one face
\[
 \sum_{j=0}^{CS}\left(1+\log^+\frac{S}{m+j}\right)\le CS;
\]
the estimate follows either by integration of $\log(S/t)$ or by summing dyadic intervals.  Irregular sides have length at most $2S$ and reflected images have bounded multiplicity.  The matrix row-sum inequality therefore bounds that variance by $CS\sum_{\partial D'}h_e^2=CP_S(D')$.  Since $B_D=B_{D'}-B_E$, we have proved
\begin{equation}\label{rg:cross-variance}
 \Var B_D\le C\{g_E+P_S(D')+T_S(D')\}.
\end{equation}

Use the squared triangle inequality only on $P_m,T_m$.  Cauchy--Schwarz, followed by the exact Gaussian moment of $B_D$, bounds $\E e^{\kappa U_m(\phi+\zeta)}$ by
\[
 (\E e^{4\kappa U_m(\zeta)})^{1/2}
 \exp\{\kappa g_D+2\kappa\delta P_m+2\kappa T_m+4\kappa^2\Var B_D\}.
\]
Subtract the desired exponent $\kappa(g_{D'}+\delta P_S+T_S)$.  By \eqref{rg:det-comparison} and \eqref{rg:cross-variance} its coefficients are bounded above by
\[
 \kappa(-1+2C\delta+4C\kappa)\,g_E,\quad
 \kappa(2\delta/L+4C\kappa-\delta)\,P_S,\quad
 \kappa(C(1+\delta)L^{-2}+4C\kappa-1)\,T_S.
\]
Choose $\delta$ so $2C\delta<1/4$, then $\kappa$ so all $8C\kappa$ terms fit within a quarter of the corresponding reserves, then $L$ large.  The same choices with reserve cover $2\kappa$.  Decrease $\kappa$ further to make the fluctuation prefactor at most $2^{|X|}$.  This proves \eqref{rg:reg-stability} with coefficient exactly $\kappa$ on the bulk energy.

For small connected $X$, its fine collar contains at most a fixed number of fine blocks and lies inside a fixed number of coarse blocks.  Equation~\eqref{rg:gradient-sup} at side $S$ gives
\[
 g_D+P_m(D)+T_m(D)\le CL^{-2}\{g_{D'}+T_S(D')\}.
\]
Now a squared triangle inequality is harmless for every term: the deterministic coefficient is $O(L^{-2})$.  The pure fluctuation bound proves the claimed half-exponent estimate for large $L$.  Finally, separated collars make the involved Gaussian subvectors independent by finite range; sums of their energies are the energy of the union, and hence their product estimates follow.  For same-scale regrouping only disjoint collars and monotonicity are needed, not Gaussian independence.
\end{proof}

\subsection{Neutral, charged, and restricted localizations}\label{rg:localization-section}

In this subsection $X$ is connected and $|X|\le n_0$, $Y=\bar X$, and $F$ obeys
\begin{equation}\label{rg:local-input}
 |F|_{m,X,\phi}\le uW_m^\kappa(X,\phi).
\end{equation}
It is periodic under a common translation of period $\omega$, and $\alpha=2\pi/\omega$.  Its charge-$q$ projection is
\[
 F_q(\phi)=\frac1\omega\int_0^\omega e^{-iq\alpha t}F(\phi+t)\,\dd t.
\]
The regulator is constant invariant; consequently every $F_q$ satisfies \eqref{rg:local-input}, and $F_q(\phi+t)=e^{iq\alpha t}F_q(\phi)$.  Let $H=\E F_0(\phi+\zeta)$.

\begin{lemma}[Neutral localization]\label{rg:localization}
Suppose the kernels satisfy Lemma~\ref{rg:kernels} and the regulator satisfies Lemma~\ref{rg:regulator}.  If $F$ is even and $X$ lies in the bulk full-block grid, put
\[
 c_X=H(0),\qquad \ell_i(v)=(v_i-v_{0,i})/m,\qquad
 D^X_{ij}=D^2H(0)[\ell_i,\ell_j].
\]
For a fine block $B\in X$ define
\begin{equation}\label{rg:quadratic-localization}
 \begin{aligned}
 \operatorname{Loc}_{X,B}H(\phi)&=\frac1{|X|}\left[c_X+
 \frac1{2m^2}\sum_{v\in B}\operatorname{ave}_{\epsilon\in\{\pm1\}^2}
 p(v,\epsilon)^TD^Xp(v,\epsilon)\right],\\
 p_i(v,\epsilon)&=m\epsilon_i\nabla_{\epsilon_i e_i}\phi(v).
 \end{aligned}
\end{equation}
Then $|c_X|+\|D^X\|\le Cu$ and
\begin{equation}\label{rg:neutral-gain}
 \left|H-\sum_{B\in X}\operatorname{Loc}_{X,B}H\right|_{S,Y,\phi}
 \le CuL^{-3}W_S^\kappa(Y,\phi).
\end{equation}
If only constants $c_X/|X|$ are subtracted, the right side is instead $CuL^{-2}W_S^\kappa$ when $F$ is even, and $CuL^{-1}W_S^\kappa$ without evenness.  Constant localization applies equally on irregular end blocks.  For a translation- and square-covariant bulk family, the sum of the quadratic pieces at each block is exactly $d e_B/2$, where $e_B=g_B$ and $|d|\le C\eps$ when $u=A^{-|X|}\eps$ and $A$ is sufficiently large.
\end{lemma}

\begin{proof}
Constant invariance permits replacing $\phi$ on the fine collar by $\phi-\phi(v_0)$.  Equation~\eqref{rg:gradient-sup} on the enclosing coarse blocks and paths of length at most $C m$ imply
\begin{align}
 \|\phi-\phi(v_0)\|_{m,D_m(X)}&\le CL^{-1}\sqrt{U_S(Y,\phi)},\label{rg:fine-radius}\\
 \|f-f(v_0)\|_{m,D_m(X)}&\le C/L,\qquad \|f\|_{S,D_S(Y)}\le1.\nonumber
\end{align}
The paths and their stencils lie in the coarse collar.  For the second differences use $m^2|\nabla^2\phi|\le CL^{-2}\sqrt{T_S}$, which is stronger than needed.  The constants depend on $R,n_0$, already fixed before $L$.

We spell out why the derivative sum and arbitrarily large real fields are compatible with Taylor expansion.  At a real interpolated field $t(\phi-\phi(v_0))$, convolution and the small-polymer part of Lemma~\ref{rg:regulator} bound the input analytic derivative sum by
\[
 C u\exp\{(\kappa/2)t^2U_S(Y,\phi)\}.
\]
After an output derivative, only the nonconstant part of its direction enters $H$.  All these directions together use at most $Ch_0/L<h_0/2$ of the fine analytic radius.  The remaining radius bounds the fixed number $r$ of Taylor derivatives by $C_rh_0^{-r}$.  Taylor's formula with its integral remainder and \eqref{rg:fine-radius} therefore bounds the output derivative sum of an order-$r$ remainder by
\begin{equation}\label{rg:taylor-bound}
 C_ruL^{-r}(1+\sqrt{U_S(Y,\phi)})^r
 e^{\kappa U_S(Y,\phi)/2}\le C_{r,\kappa}uL^{-r}W_S^\kappa(Y,\phi).
\end{equation}
The last inequality is the finite supremum of $(1+\sqrt t)^re^{-\kappa t/2}$ over $t\ge0$.  This argument justifies the derivative sums by their nonnegative absolute majorants.  Evenness removes the linear Taylor term.  Thus constant subtraction gives order two in the even case and order one otherwise.

For quadratic localization, Taylor expansion to order two leaves the order-three bound \eqref{rg:taylor-bound}.  For each $v\in B$ and $\epsilon$, replace the field modulo constants by the affine function with coordinate slopes $\epsilon_i\nabla_{\epsilon_i e_i}\phi(v)$.  The error in fine test norm is at most $CL^{-2}\sqrt{U_S}$, by summing second differences along coordinate paths.  On output directions its norm is at most $CL^{-2}$.  The difference between the two quadratic forms contains one such error and one factor of size $CL^{-1}(1+\sqrt{U_S})$, or two errors.  Applying the same radius allocation and exponent slack gives $CuL^{-3}W_S^\kappa$.  Every full fine block has $m^2$ sites; averaging those affine quadratic forms over $v,\epsilon$ and then over $B\in X$ gives exactly \eqref{rg:quadratic-localization}.  At zero field the derivative bound and $\|\ell_i\|_{m,D_m(X)}\le C$ prove $|c_X|+\|D^X\|\le Cu$.

Finally define $D_B=\sum_{X\ni B}D^X/|X|$.  Translation covariance makes it independent of $B$; reflection changes the sign of each off-diagonal entry, and a quarter-turn interchanges the diagonal entries.  Hence $D_B=d\Id$.  Direct substitution, retaining both signs of each lattice direction, yields
\[
 \frac1{2m^2}\sum_{v\in B}\operatorname{ave}_{\epsilon}p^T(d\Id)p
 =\frac d4\sum_{v\in B}\sum_{e=\pm e_i}(\nabla_e\phi(v))^2
 =\frac d2 e_B(\phi).
\]
The rooted animal count in Lemma~\ref{rg:geometry} bounds $|d|$ by $C\eps\sum_{n\le n_0}(C_R/A)^n\le C\eps$.  Quadratic localization was asserted only for full bulk blocks; no replacement of an end-block volume by $m^2$ has been made.
\end{proof}

\begin{lemma}[Charged contraction]\label{rg:charged}
There are choices of $h_0$, then $L$, then a lower bound $\alpha_0$, such that for $\alpha\ge\alpha_0$ the sum of all nonzero charges satisfies
\begin{equation}\label{rg:charge-gain}
 \sum_{q\ne0}|\E F_q(\phi+\zeta)|_{S,Y,\phi}
 \le CuL^{-4}W_S^\kappa(Y,\phi).
\end{equation}
The same argument with shift amplitude $a>0$ has charge gain
\begin{equation}\label{rg:general-charge-gain}
 \exp\{-(a|q|\alpha-a^2/2)\Gamma_m(v_0,v_0)+h_0|q|\alpha\},
\end{equation}
provided $h_0$ is larger than a fixed constant times $a$.  It applies whenever this gain is summable.  Fundamental-charge localization at marginal frequency requires a separate subtraction and is not included in \eqref{rg:charge-gain}.
\end{lemma}

\begin{proof}
Set $w(v)=\Gamma_m(v,v_0)$, $w_0=w(v_0)$.  The binary positive-order kernel estimates and $\diam D_m(X)\le Cm$ imply
\[
 \|w-w_0\|_{m,D_m(X)}\le C,
\]
uniformly in $L$.  The Gaussian shift identity, with $\tau=\operatorname{sgn}(q)$, is
\begin{equation}\label{rg:imaginary-shift}
 \E F_q(\phi+\zeta)
 =e^{a^2w_0/2}\E\left[e^{-ia\tau\zeta(v_0)}
 F_q(\phi+\zeta+ia\tau w)\right].
\end{equation}
The shift $ia\tau w_0$ is constant and contributes $e^{-a|q|\alpha w_0}$.  Its nonconstant remainder has fine norm at most $Ca$.  A coarse direction is its anchor value, of absolute value at most one, plus a nonconstant part of fine norm at most $C/L$.  The first contributes $|q|\alpha$ per derivative; summing it gives $e^{h_0|q|\alpha}$.  Choose $h_0>4Ca$ and $L$ so large that the imaginary remainder and all residual directions consume less than $h_0$.  The fine derivative sum then bounds their full contribution by the input norm at the real field $\phi+\zeta$.  Lemma~\ref{rg:regulator} gives \eqref{rg:general-charge-gain} times $CuW_S^{\kappa/2}$.

For rigor in \eqref{rg:imaginary-shift}, use the Gaussian eigen-coordinates with positive eigenvalues.  Their exact real Cameron--Martin identity is $\E F_q(\phi+\zeta)=\E[e^{-z\zeta(v_0)-z^2w_0/2}F_q(\phi+\zeta+zw)]$ for real $z$.  Both sides continue analytically in the scalar shift in a neighborhood containing the segment $[0,ia\tau]$: the constant character handles its constant part and the stated radius bound handles its residual part.  A slightly larger Gaussian regulator parameter bounds the integrands and their differentiated series along this compact segment.  Dominated continuation proves the identity.  Null covariance directions do not enter $w$ and cause no change.

Take $a=1$.  The diagonal lower bound $w_0\ge c\log L$ allows first choosing $L$ so $w_0\ge2h_0+2$ and then $\alpha_0$ so
\[
 (\alpha|q|-1/2)w_0-h_0\alpha|q|\ge4\log L+c'\alpha(|q|-1)\log L.
\]
The resulting geometric series over $|q|\ge1$ proves \eqref{rg:charge-gain}.  Fourier projection and inversion are valid for values and each derivative because periodic analyticity gives exponentially decaying Fourier coefficients on every compact real-field neighborhood.  The obtained summable output-norm majorant permits summing the full analytic derivative series as well.
\end{proof}

\begin{corollary}[Placement factors]\label{rg:restricted-localization}
For small linear terms at a prescribed coarse block the number of fine placements is at most $CL^2$.  For terms whose fine support is within a fixed number of fine blocks of a specified straight side it is at most $CL$, and for terms within a fixed number of fine blocks of one of a fixed number $k$ of marked points it is at most $Ck$.  Thus the neutral bulk, even wall, and unrestricted marked remainders of Lemma~\ref{rg:localization}, after summing placements, are bounded respectively by
\[
 C L^{-1}\eps,\qquad C L^{-1}\eps,\qquad C kL^{-1}\eps
\]
in output activity norm, apart from large polymers and nonlinear terms.  A bulk term whose fluctuation neighborhood touches a wall can always be bounded without subtraction by $C_{L,A}\eps$ and used as forcing.
\end{corollary}

\begin{proof}
A small support contains at most $n_0$ blocks, with boundedly many shapes at its anchor.  A coarse rectangle has at most $CL^2$ fine anchors; its fixed-width boundary strip has at most $CL$; a fixed-width neighborhood of a point has at most $C$.  End rectangles have side between $S$ and $2S$, leaving these counts unchanged.  Since $|\bar X|\le|X|$, the activity weight ratio is at most one.  Multiply the counts by $L^{-3},L^{-2},L^{-1}$, respectively.  The last assertion follows directly from convolution, the regulator bound, and the finite number $C_{L,A}$ of small supports and assigned pieces.  It invokes no bulk cancellation across a physical wall.
\end{proof}

\subsection{Exact integration, relocation, and constant extraction}

For a finite block grid $\mathcal B_m$, define the subset partition functional
\begin{equation}\label{rg:partition}
 \mathcal Z_m(K;\phi)=\sum_{V\subset\mathcal B_m}
 \prod_{X\in\operatorname{Comp}_R(V)}K(X,\phi),
\end{equation}
with empty product one.  Equivalently sum over collections of nonempty $R$-connected polymers whose mutual index distances exceed $R$.  This is a finite algebraic sum; no logarithm or infinite-volume partition function is required to define the local map.

\begin{proposition}[Exact local partition map]\label{rg:partition-map}
Under the locality, finite-range, and norm hypotheses already specified, shell integration, localizations on small linear inputs, and extraction of their designated field-independent coefficients define activities $K'$ and numbers $a_b$, one at each coarse block, satisfying the exact identity
\begin{equation}\label{rg:exact-map}
 \E\mathcal Z_m(K;\phi+\zeta)
 =\prod_{b\in\mathcal B_S}(1+a_b)\mathcal Z_S(K';\phi).
\end{equation}
The map preserves the period, locality, and every symmetry respected by the designated localizations.  To first order, $K'$ is the sum of the subtracted small expectations, the unchanged large expectations, and the nonconstant parts assigned to coarse singletons.  When singleton polynomials are measured separately, write
\[
 K(X)=K_{\rm reg}(X)+\ind_{X=\{B\}}\sum_{\nu\in I}t_{B,\nu}P^\nu_{m,B},
 \qquad
 |P^\nu_{m,B}|_{m,\{B\},\phi}\le C_{\rm pol}W_m^\kappa(\{B\},\phi).
\]
Here $I$ is fixed and finite, and the $P^\nu_{m,B}$ form a fixed family of local polynomials, independent of the activity inputs, with the prescribed collar and period.  The displayed point-norm bound is uniform in $\nu$, the retained scale, the block, and the finite volume.  Suppose
\[
 \|K_{\rm reg}\|_{m,A,\kappa}
       +\sup_B\sum_{\nu\in I}|t_{B,\nu}|\le\eps.
\]
If a fixed scalar representation retains $K_{\rm reg}=\sum_\mu K_{\rm reg}^\mu$ as separate components, replace $\|K_{\rm reg}\|_{m,A,\kappa}$ in this bound by $\sum_\mu\|K_{\rm reg}^\mu\|_{m,A,\kappa}$, and include every separately retained coefficient occurrence in the absolute coefficient sum.  With external parameters fixed, all designated localization and scalar extraction maps are fixed complex-linear maps of this measured input vector.  Let $F_{Y,b}$ be the complete singleton piece assigned from the small linear inputs with original closure $Y$ to $b\in Y$, retaining $Y$ as its exclusion label, and let $a_b$ be the sum of the designated field-independent scalar contributions from these same local $Y\to b$ inputs.  Each assigned piece has the prescribed target collar and period.  For every complex input vector obeying the displayed bound, require
\[
 \sup_b\left\{
 \sum_{Y\ni b}\sup_\phi
 \frac{|F_{Y,b}|_{S,\{b\},\phi}}{W_S^\kappa(\{b\},\phi)}
       +|a_b|\right\}\le C_{L,A}\eps,
\]
uniformly in retained scales and finite volumes.
Then on a sufficiently small ball
\begin{equation}\label{rg:nonlinear-map}
 \|K'-K'_{\rm lin}\|_{S,A,\kappa}\le C_{L,A}\eps^2.
\end{equation}
This is a complex analytic map of its activity inputs, with the corresponding convergent Taylor series and derivative bounds.  The assertion is not analyticity in the frequency convention $\alpha$.

The same finite algebra and spatial enumeration apply to a separately specified representation under the same locality, finite-range, and exclusion hypotheses, with fixed complex-linear designated maps, provided its input size includes every retained component and it supplies the corresponding bounds for compatible-family products, Gaussian integration, complete assigned singleton pieces, and scalars in its normalized size.  The point-norm and complex analyticity assertions above are statements about the scalar representation.

After $L$ has been fixed, taking $A$ sufficiently large makes the contribution from large linear inputs an arbitrarily small multiple of their norm.  The small linear estimates are those of Lemmas~\ref{rg:localization}, \ref{rg:charged} and Corollary~\ref{rg:restricted-localization}.  All constants and smallness radii are uniform in finite volumes.  For the terminal assertion, separately assume that the retained activity grid at scale $m_*$ satisfies the grid conditions above and $M/m_*\le D_0L$; its block count is then at most $C(D_0L)^2$.  The spectral decomposition and terminal estimates retain their own stated domains.
\end{proposition}

\begin{proof}
We give the algebra before bounding it.  For a configuration of compatible fine polymers take the union of their coarse closures and its $r_0$-connected components.  Fine fields belonging to different such components have mutual distance greater than the shell range: the support collar uses only a fixed number of lattice steps and $r_0>4p+10$.  Their Gaussian subvectors are independent.  For each temporary component $Y$ sum the expectations of products of all compatible fine families with closure union exactly $Y$; call this sum $\widetilde K(Y)$.  Independence proves that the integrated partition functional is precisely the temporary-radius partition expansion in these $\widetilde K$.  Its linear term is $\sum_{\bar X=Y}\E K(X)$.

Apply the fixed designated localization to each small linear summand, using Lemma~\ref{rg:localization} for the ordinary choice.  Collect pieces assigned to a coarse block $b\in Y$ as $F_{Y,b}$.  Write identically
\[
 \widetilde K(Y)=\left(\widetilde K(Y)-\sum_{b\in Y}F_{Y,b}\right)+\sum_{b\in Y}F_{Y,b}.
\]
Expand this choice independently in every temporary configuration.  A residual retains support $Y$; a chosen piece receives support $\{b\}$, while its original exclusion information $Y$ is retained.  Group the union of these new supports into its $R$-connected components.  This factorizes even with the retained exclusions.  Indeed a support that shrinks came from a small linear input and has index diameter at most one.  If two original supports violate a temporary exclusion, their new supports have distance at most $r_0+2<R$.  They therefore belong to the same final component.  There is no exclusion constraint between distinct final components.  Within a component it is retained in the defining finite sum.  Denote the resulting final-radius activities by $H(Y)$.

Let $a_b$ be the sum of the designated field-independent coefficients assigned at $b$.  Assume $|a_b|<1/2$ and put $d_b=(1+a_b)^{-1}-1$.  Multiply the partition expansion with activities $H$ by
$\prod_b(1+d_b)$.  View a chosen $d_b$ as a constant decoration that may overlap an occupied old block.  Group the occupied union again by radius $R$.  Since the scalar multiplier is exactly $\prod_b(1+a_b)^{-1}$, this defines $K'$ and proves \eqref{rg:exact-map}.  At first order $d_b=-a_b$, so precisely the designated constants cancel.  A translated gradient polynomial can have a nonzero value at zero field; if that value is not designated for extraction, the algebra keeps it in $K'$.  This observation will matter for defects.

We next justify the norm bounds, including uniformity of the sums.  Initially use the stronger output weight $4A$.  Expand each singleton input into its part measured by the activity norm and its separately measured polynomial parts.  After summing the absolute bounds of any separately retained regular components, an ordinary input has point norm at most $\eps A^{-|X|}W_m^\kappa$.  After summing $\nu$, a separately measured polynomial singleton has point norm at most $C_{\rm pol}\sum_\nu|t_{B,\nu}|W_m^\kappa\le(C_A\eps)A^{-1}W_m^\kappa$, with $C_A=AC_{\rm pol}$.  For $P=e_B/2$, the uniform local bound follows from $e_B=g_B\le U_m$ and the bounds $C\sqrt{U_m}$ and $C$ on its first and second derivatives in unit test directions; the regulator absorbs these terms.  Compatible fine polymers have disjoint collars.  The product inequality and Lemma~\ref{rg:regulator} give, for a family of total size $n$ and $j$ components, of which $j_{\rm pol}$ are separately measured polynomial singletons after summing their indices, a bound
\[
 \eps^j C_A^{j_{\rm pol}}(2/A)^n W_S^\kappa(Y).
\]
For each large component $X_i$, \eqref{rg:closure} gives $|\bar X_i|\le |X_i|/2$, whereas every component obeys $|\bar X_i|\le|X_i|$.  Hence multiplication by $(4A)^{|Y|}$ costs at most a product with factor $8^{|X_i|}A^{-|X_i|/2}$ on a large component and $8^{|X_i|}$ on a small one.  Use \eqref{rg:family-count}.  For fixed $L$ the sum over the size of each large component is bounded by
\[
 \sum_{n>n_0}(C A^{-1/2})^n,
\]
and tends to zero as $A\to\infty$.  The sum over each small component has at most $n_0$ terms and a finite cost depending on the geometry, with the additional fixed $C_A$ cost only for separately measured singletons.  Large components retain their original weight gain.  The long-jump factor $(CL^2)^{2j}$ is a fixed factor per component.  Thus all families with $j\ge2$ sum to at most
\[
 \sum_{j\ge2}(C_{L,A}\eps)^j\le C_{L,A}\eps^2
\]
when $\eps<(2C_{L,A})^{-1}$.  The $j=1$ terms are bounded linearly, and the large ones can be made an arbitrarily small linear multiple by increasing $A$.  Bounding the rooted sum over all output supports is stronger than the required supremum over one output support, so this argument proves the claimed norm estimate without a volume factor.

For the ordinary localization, $|c_X|+\|D^X\|\le Cu$ and the normalized formula \eqref{rg:quadratic-localization} give the required complete singleton bound: its value is at most $Cu(1+U_S)$, and its first two derivatives in unit coarse directions are at most $Cu\sqrt{U_S}$ and $Cu$.  The regulator absorbs these terms.  The finite small-support count then verifies the stated rooted assignment and scalar bound; for other designated maps use that bound directly.  In the relocation estimate keep each $F_{Y,b}$ as one atom carrying its original label $Y$, with support $Y$ for its subtraction and $\{b\}$ for its assignment.  Its fine summands remain combined.  A small $Y$ has index diameter at most one, so its possible labels at a fixed $b$ and the weight change between these two supports have a fixed $C_{L,A}$ cost; regulator monotonicity handles the support change.  Distinct temporary objects have distance greater than $r_0$, and shrinking a small label by at most one block leaves their assigned collars disjoint.  Count a final connected family by the previous tree of fine objects with one added coarse vertex for each such atom.  An added vertex has only a fixed number of label choices, and its tree attachments have at most $CL^2$ fine positions or a fixed number of coarse positions.  Marking attachment positions costs only an exponential factor in the old fine size and the number of added atoms.  The rooted assignment bound supplies $C_{L,A}\eps$ per added vertex.  Thus the same exponential family estimate applies, with large objects retaining their original weight gain.  The bound at weight $4A$ remains available after relocation, and terms involving at least two objects are bounded by the same geometric series starting at degree two.  The subtraction and designated singleton insertion are complex-linear, so all other changes so far are included in \eqref{rg:nonlinear-map}.

For clarity, the last decoration sum can be bounded directly without discarding overlaps.  Fix a final support $Y$ and its old occupied set $V\subset Y$.  The old components are determined by $V$.  Decorations are compulsory on $Y\setminus V$ and optional on $V$.  Put $d_* =\sup_b|d_b|\le C_{L,A}\eps$.  Monotonicity of the regulator and disjointness of old collars imply that the absolute sum for this $V$ is at most its old activity coefficient times
\[
 d_*^{|Y\setminus V|}(1+d_*)^{|V|}W_S^\kappa(Y).
\]
The old weight $4A$, compared with the final weight $A$, supplies $4^{-|V|}$; optional overlaps cost at most $2^{|V|}$.  The holes supply $(Ad_*)^{|Y\setminus V|}$.  Thus the subset sum is bounded by
\[
 \sum_{V\subset Y}2^{-|V|}(Ad_*)^{|Y\setminus V|}
 =(1/2+Ad_*)^{|Y|}\le(5/8)^{|Y|}
\]
when $Ad_*\le1/8$, leaving an exponential reserve.  A term with an old activity and a hole contains at least one factor $\eps$ from each.  If there are no holes, the nonempty-overlap sum is bounded by $|Y|d_*(1+d_*)^{|Y|}$; the exponential reserve absorbs $|Y|$.  Pure decorations on a support of at least two blocks are degree at least two; on a singleton, $d_b=-a_b+O(a_b^2)$.  These bounds give the stated quadratic estimate at the last step.

The designated maps are bounded complex-linear maps of the measured input vector.  All remaining operations are sums of bounded multilinear maps and the absolutely convergent scalar inverse series for $1+a_b$.  Their absolute bounds on a slightly larger input ball imply convergence of the complex Taylor series and its termwise derivatives on the smaller ball.  Symmetry and periodicity are preserved by the explicit operations.  Locality is preserved because every field used by an output object is inside its output support and fixed collar; its numerical coefficients can depend on predecessor data only in a fixed larger neighborhood.  Iteration expands that neighborhood by the geometric sum of earlier scale lengths, bounded by a constant times the present scale.  Choosing $D_0$ larger than this constant ensures exact plane/torus agreement on unwrapped neighborhoods through the stopping scale.  Finally, under the hypothesis on the terminal activity grid, $M/m_*\le D_0L$ bounds the number of terminal blocks.  When the terminal covariance also satisfies \eqref{rg:terminal-reg}, that estimate controls their final Gaussian integrations.  Under these terminal hypotheses, if every terminal activity norm tends to zero, then its finite subset expansion tends to one after the designated scalar factors have been removed.
\end{proof}

The parameter order used below is now explicit: fix interaction and connection ranges and $n_0$; choose $\delta$, then $\kappa$ (with the required exponent and H\"older reserves); choose $h_0$ large enough for the specified imaginary translations; choose $L$ large enough for the geometric and small-support gains; choose $A$ large enough for large supports; finally choose the activity ball small enough for the inverse series and nonlinear estimates.  A large-frequency construction chooses its frequency lower bound after $h_0,L$.  An entry construction using an additional regulator must prove its Gaussian composition and determinant bounds before applying Proposition~\ref{rg:partition-map}.
